\section{System Overview and Design}
\label{sec:design}
Figure~\ref{fig:system} separates Application prediction, Router execution
ownership and Store-Master policy. Workers seal and copy primary KV; the
standby holds backup and recovery allocations separately. Control messages
carry metadata, while workers and Store exchange KV bytes.

\begin{figure*}[tp]
\centering\includegraphics[width=\textwidth]{figures/system-overview.pdf}
\caption{AegisKV architecture. Solid blue and orange arrows are independent
source-owned protection paths; dashed purple arrows carry control metadata.
HBM and DRAM have separate reservations and acknowledgements. The Controller is inside the native Store Master. Backup HBM and recovery
primary KV occupy distinct allocations. The Application
uses GBDT in the prototype; simulation uses the explicitly stated ideal-duration
CPM projection. The study excludes SSD.}
\Description{Application events feed Slack prediction and tickets; a Router,
Store Master and quota controller direct a source worker, a route-bound HBM
backup target and an independent DRAM store.}
\label{fig:system}
\end{figure*}

\subsection{Failure model and compatibility}
AegisKV targets fail-stop loss of worker KV with a surviving Application,
Router and Controller. Recovery requires a live replica matching model,
KV ABI, dtype and parallel layout, with all shards for each logical chunk.
The evaluated DP configurations use TP=PP=1; partitions, control-plane
failure and cross-layout recovery are outside scope.

Rendezvous hashing of request identity, seed, role and worker binds a primary
and compatible anti-colocated standby. Standby failure rotates the backup;
source failure fences and versions the route before activating that standby.
Remote HBM cannot be relabeled as local recovery state.

\subsection{Algorithm 1: causal residual-workflow Slack}
On a request or tool return, the Application updates the observed execution
summary and constructs the causally visible residual DAG to the next join.
A GBDT predicts remaining duration for active agent nodes and full duration
for known, unstarted nodes; structural nodes have zero weight. Standard
forward and backward CPM compute earliest and latest starts. For active
branch $i$, raw float is $s_i^{raw}=\pos{LS_i-ES_i}$; a calibrated margin
$m$ gives the published signal $s_i=\pos{s_i^{raw}-m}$.
Algorithm~\ref{alg:slack-update} gives the complete update.

Unknown branches are excluded. Incomplete structure or missing predictions
invalidate all active sibling tickets, setting Slack to zero. Calibration
uses workflow-disjoint data. The Controller ages tickets as
$s_i(t)=\pos{s_i-(t-t_i)}$ until replacement by a valid update.
Predicted CPM float does not guarantee future Slack. Graph processing costs
$O(|V|+|E|)$ with a fixed predictor.

% Source: Mooncake application/algorithm1.py:update_slack, inspected 2026-09-06.
\begin{algorithm}[t]
\caption{Causal GBDT--CPM update with calibrated Slack}
\label{alg:slack-update}
\begin{tabular}{@{}r@{\hspace{0.5em}}>{\raggedright\arraybackslash}p{\dimexpr\columnwidth-2.5em\relax}@{}}
\multicolumn{2}{@{}>{\raggedright\arraybackslash}p{\columnwidth}@{}}{\textit{Input:} event $e$, time $t$, causal predictor $f_\theta$, workflow-disjoint margin $m\ge0$. $A$: active nodes; $v$: event version.}\\
\multicolumn{2}{@{}>{\raggedright\arraybackslash}p{\columnwidth}@{}}{\textbf{UpdateSlack}$(e,t)$}\\
1 & $(G=(V,E),A,\phi,v)\gets\operatorname{Observe}(e,t)$; $\pi\gets\operatorname{Topo}(G)$.\\
2 & \textbf{if} graph or node metadata is incomplete: publish invalid zero tickets for all $i\in A$; \textbf{return}.\\
3 & For each $u\in V$, set $w_u\gets f_\theta(\phi_u)$ for agent nodes and $w_u\gets0$ for control nodes.\\
4 & \textbf{if} any prediction is unavailable: publish invalid zero tickets for all $i\in A$; \textbf{return}.\\
5 & For $u\in\pi$, set $ES_u\gets0$ at roots, otherwise $ES_u\gets\max_{z\in\operatorname{pred}(u)}(ES_z+w_z)$.\\
6 & $C\gets\max_{u:\operatorname{succ}(u)=\varnothing}(ES_u+w_u)$.\\
7 & For $u\in\operatorname{rev}(\pi)$, set $LS_u\gets C-w_u$ at sinks, otherwise $LS_u\gets\min_{z\in\operatorname{succ}(u)}LS_z-w_u$.\\
8 & For each $i\in A$, set $s_i^{raw}\gets[LS_i-ES_i]^+$ and $s_i\gets[s_i^{raw}-m]^+$.\\
9 & Publish valid tickets $(i,s_i,s_i^{raw},t,v)$ with prediction and calibration provenance.\\
\end{tabular}
\end{algorithm}


\subsection{Algorithm 2: independent tier-local admission}
The worker seals 4,096-token prompt and 256-token decode chunks. Final prompt
chunks may be shorter; open decode tails remain unprotected. Model/layout
bytes determine allocation. Each sealed chunk independently enters admission
for Router-bound HBM and direct Store DRAM.

Let $f_i^r$ be request $i$'s committed prefix in tier $r$. The tier-local
recovery estimate includes loading that prefix, recomputing the uncovered
prompt and silently replaying unprotected output positions. It is capped
by full recomputation. Define
\begin{equation}
J_i^r(t)=\pos{T_i^r(f_i^r,t)-s_i(t)}.
\end{equation}
For candidate $x$ and tentative victim set $Y$, the gain and joint loss are
\begin{align}
A_x &= \pos{J_i^r(\mathrm{before})-J_i^r(\mathrm{after}\ x)},\\
D_Y &= \sum_{j\in\operatorname{req}(Y)}
       \pos{J_j^r(\mathrm{after}\ Y)-J_j^r(\mathrm{before})}.
\end{align}
Loss is recomputed once per affected request from the resulting prefix;
summing independent per-chunk deletion costs could double-count recovery.
The other tier cannot mask $A_x$ or $D_Y$. Backup copy time is an actual
runtime cost, not a speculative future destaging deduction from HBM gain.

Each pool $r$ retains rent $L_r$ and one exposed, safely removable suffix tail
per request. Chunk $c$ has key $H_c=L_r+A_c/b_c$ at commitment. Slack or
frontier changes refresh exposed deletion-loss keys while preserving credit:
$H_c\gets\max\{L_r,H_c+(D_c^{new}-D_c^{old})/b_c\}$. Hidden predecessors keep
their saved keys until exposed. This avoids issuing fresh credit merely
because rent advanced.

Algorithm~\ref{alg:slack-landlord} selects tentative tails with keys strictly
below $\Theta_x=L_r+A_x/b_x$, excluding the candidate's own request and unsafe
or read-pinned data. It exposes predecessors only in suffix order, stopping
when enough real bytes have been released. Admission requires the strict
inequality $A_x>D_Y$; zero gain and ties bypass. A successful decision evicts
$Y$, advances rent, reserves final destination space and starts an independent
source copy. The two tiers may produce HBM-only, DRAM-only, both, or neither.
No HBM-to-DRAM destaging operation exists.

% Native independent-tier policy; implementation details are discussed in text.
\begin{algorithm*}[tp]
\caption{\textsc{SlackLandlord}: independent admission and asynchronous tier commitment}
\label{alg:slack-landlord}
\textit{Symbols:} $x=(c,r)$: sealed chunk $c$ in pool $r$; $i=\operatorname{req}(c)$; $b_c$: bytes;
$L_r$: rent; $Q_r$: heap exposing one safe suffix tail per request; $Y$: tentative victims.
\par\medskip
\noindent\begin{minipage}[t]{.48\textwidth}
\vspace{0pt}
\begin{tabular}[t]{@{}r@{\hspace{0.5em}}>{\raggedright\arraybackslash}p{\dimexpr\linewidth-2.5em\relax}@{}}
\multicolumn{2}{@{}>{\raggedright\arraybackslash}p{\columnwidth}@{}}{\textbf{SlackLandlord}$(c,t)$}\\
1 & $x_H\gets\operatorname{BuildHbm}(c)$; $x_D\gets\operatorname{BuildDirectDram}(c)$.\\
2 & $h\gets\operatorname{TryAdmit}(x_H,t)$.\\
3 & $d\gets\operatorname{TryAdmit}(x_D,t)$.\\
4 & $\operatorname{StartAcceptedSourceCopies}(h,d)$.\\
5 & $\operatorname{return}\ \operatorname{ResidencyOutcome}(h,d)$.\\
\multicolumn{2}{@{}>{\raggedright\arraybackslash}p{\columnwidth}@{}}{\textbf{MinTail}$(Q,Y,i,t)$}\\
6 & \textbf{while} $Q\ne\varnothing$ \textbf{do}\\
7 & $\quad(y,H_y)\gets Q.\operatorname{min}()$.\\
8 & $\quad$\textbf{if} $\operatorname{Stale}(y,Y,t)$ \textbf{then} $Q.\operatorname{update}(y,\operatorname{RefreshKey}(y,Y,t))$; \textbf{continue}.\\
9 & $\quad Q.\operatorname{pop}()$; $\operatorname{req}(y)\ne i\Rightarrow\operatorname{return}\ (y,H_y)$.\\
10 & \textbf{end while}; $\operatorname{return}\ (\bot,\infty)$.\\
\multicolumn{2}{@{}>{\raggedright\arraybackslash}p{\columnwidth}@{}}{\textbf{OnCopyDone}$(x=(c,r),t')$}\\
11 & \textbf{if} $\operatorname{AllAck}(x)\land\operatorname{Valid}(x)$ \textbf{then}\\
12 & $\quad H_c\gets L_r+\operatorname{Gain}(x,t')/b_c$; $\operatorname{state}(x)\gets\textsc{Committed}$; $Q_r.\operatorname{insert}(c,H_c)$; \textbf{return}.\\
13 & $\operatorname{state}(x)\gets\textsc{Aborted}$; $\operatorname{Free}(c,r)$.\\
\end{tabular}
\end{minipage}\hfill
\begin{minipage}[t]{.48\textwidth}
\vspace{0pt}
\begin{tabular}[t]{@{}r@{\hspace{0.5em}}>{\raggedright\arraybackslash}p{\dimexpr\linewidth-2.5em\relax}@{}}
\multicolumn{2}{@{}>{\raggedright\arraybackslash}p{\columnwidth}@{}}{\textbf{TryAdmit}$(x=(c,r),t)$}\\
14 & $A\gets\operatorname{Gain}(x,t)$; $\Theta\gets L_r+A/b_c$; $Y\gets\varnothing$; $Q\gets\operatorname{View}(Q_r)$.\\
15 & $f\gets\operatorname{Free}(r)$.\\
16 & \textbf{while} $f<b_c$ \textbf{do}\\
17 & $\quad(y,H_y)\gets\operatorname{MinTail}(Q,Y,i,t)$.\\
18 & $\quad$\textbf{if} $(y=\bot)\lor(H_y\ge\Theta)$ \textbf{then} \textbf{return} \textsc{Bypass}.\\
19 & $\quad Y\gets Y\cup\{y\}$; $f\gets f+b_y$.\\
20 & $\quad z\gets\operatorname{PrevTail}(y,Y)$; $z\ne\bot\Rightarrow Q.\operatorname{insert}(z,\operatorname{RefreshKey}(z,Y,t))$.\\
21 & \textbf{end while}\\
22 & $D\gets\operatorname{Loss}(Y,t)$; $A\le D\Rightarrow\operatorname{return}\ \textsc{Bypass}$.\\
23 & $\operatorname{Evict}(Y)$; $Y\ne\varnothing\Rightarrow L_r\gets\max\{L_r,\max_{y\in Y}H_y\}$.\\
24 & $\operatorname{Allocate}(c,r,b_c)$; $\operatorname{state}(x)\gets\textsc{Pending}$.\\
25 & $\operatorname{owner}(x)\gets\operatorname{SourceWorker}(c)$.\\
26 & $\operatorname{return}\ \operatorname{TierDecision}(r)$.\\
\end{tabular}
\end{minipage}
\end{algorithm*}


\subsection{Commit, replacement and request lifetime}
Each admitted path starts PENDING and becomes COMMITTED only after all shards
transfer under a valid route and layout. Failure aborts that path and releases
its reservation. Figure~\ref{fig:commit} shows divergent frontiers when failure
falls between the tiers' acknowledgements.

The connector advances the logical ordinal after its enabled paths resolve.
When one tier commits and the other bypasses, the latter seals its lagging
prefix. All-bypass outcomes permit retrying the same frontier. Admission
remains independent; neither tier requires the other's residency.

\begin{figure*}[tp]
\centering\includegraphics[width=\textwidth]{figures/commit-recovery.pdf}
\caption{Commit visibility and physical recovery. (a) Source failure occurs
after HBM commits chunk 1 and before DRAM acknowledges it. (b) HBM retains
chunks 0--1; DRAM retains only chunk 0. C, X and dash denote committed, aborted
and absent state. (c) Recovery copies the surviving HBM prefix into separately
allocated primary KV and recomputes suffix 2. Event spacing is schematic;
no duration or pressure-race guarantee is inferred.}
\Description{A logical state timeline, per-tier committed-prefix arrays,
and a backup-to-primary memory copy preceding suffix recomputation.}
\label{fig:commit}
\end{figure*}

Replacement deliberately has a protection gap: victims disappear before the
new copy commits. The policy does not stage both generations or restore
victims after a transfer failure. Recovery therefore uses only the surviving
committed prefix at the actual fault instant. When a request finishes,
CompleteRequest drains pending paths before reclaiming both tiers' metadata,
Store keys and HBM slots. Per-request cache salts isolate content ownership
until shared physical reference counting is supported.

\subsection{Algorithm 3: HBM quota and foreground pressure}
For registered HBM capacity $M_p$, target $p$ receives quota
$\lfloor\alpha_pM_p\rfloor$. This envelope is separate from the physical
allocator's accounting for weights, primary KV and runtime memory. A pressure
observation means insufficient primary headroom or an allocation/reclaim
stall. Algorithm~\ref{alg:guarded-aimd} multiplies $\alpha_p$ by $\gamma$ on
pressure and adds $\delta$ after $K$ stable windows, within configured bounds.

The local enforcer blocks new protection when occupancy exceeds the quota,
and reclaims safe HBM tails until capacity permits admission. HBM-only tails
may be removed; a DRAM copy is not a prerequisite. The design requires a stable read source while materializing recovery.
The simulator pins selected sources; the prototype uses local copy locking
and recovery visibility leases, with a remaining pre-load pressure race
(Section~\ref{sec:implementation}). AIMD selects no
chunks and schedules no transfers. Transfer Engine retains ownership of copy
queues and bandwidth. The control update costs $O(P)$ for $P$ targets.

% Compiled by sections/03-system-design-v2.tex.
% Design-SHA256: 0fdfd097a536ab73cd9f7eb12099eb6f5b9c068a061da30af57a8e6798caacac
\begin{algorithm}[tp]
\caption{\textsc{HBMQuotaAIMD}: Peer-HBM protection share}
\label{alg:guarded-aimd}
\begin{tabular}{@{}r@{\hspace{0.5em}}>{\raggedright\arraybackslash}p{\dimexpr\columnwidth-2.5em\relax}@{}}
\multicolumn{2}{@{}>{\raggedright\arraybackslash}p{\columnwidth}@{}}{\textit{Symbols.} $p$: Peer-HBM target; $\alpha_p$: protection fraction; $q_p\in\{0,1\}$: HBM pressure; $z_p$: stable windows.}\\
\multicolumn{2}{@{}>{\raggedright\arraybackslash}p{\columnwidth}@{}}{$0<\gamma<1$: decrease factor; $\delta>0$: increase step; $K$: stable-window threshold; $[\alpha_p^{min},\alpha_p^{max}]$: bounds.}\\
\multicolumn{2}{@{}>{\raggedright\arraybackslash}p{\columnwidth}@{}}{\textbf{ControlWindow}$(t)$}\\
1 & \textbf{for} $p\in\mathcal P$ \textbf{do}\\
2 & $\quad$\textbf{if} $q_p=1$ \textbf{then}\\
3 & $\qquad\alpha_p\gets\max\{\alpha_p^{min},\gamma\alpha_p\}$; $z_p\gets0$.\\
4 & $\quad$\textbf{else} $z_p\gets z_p+1$.\\
5 & $\qquad z_p\ge K\Rightarrow(\alpha_p,z_p)\gets(\min\{\alpha_p^{max},\alpha_p+\delta\},0)$.\\
6 & $\quad$\textbf{end if}; $\operatorname{Publish}(p,\alpha_p,t)$.\\
7 & $\quad\operatorname{Storage.EnforceHbmQuota}(p,\alpha_p)$; \textbf{end for}.\\
\end{tabular}
\end{algorithm}


\subsection{Recovery and invariants}
The Controller orders ready recovery requests by increasing remaining Slack,
then decreasing current recovery harm and request identity. The activated
worker loads target-local HBM first and uses DRAM for the uncovered protected
suffix. These copies populate separate primary KV buffers: local HBM is not
free zero-copy activation. Before taking read leases, primary capacity may
reclaim backup tails and expose DRAM or recomputation fallback. If other live
primary allocations prevent admission, recovery waits on the bound worker.

The executor recomputes the uncovered prompt, silently replays to the Router's
committed output position, then resumes visible output. Pending data never
contributes to restored tokens. Per-tier prefix continuity, version checks,
capacity conservation and output fencing remain mandatory regardless of
Slack, rent or quota. The design accepts reduced redundancy during
replacement, and makes no global-optimum or competitive-ratio claim for the
coupled multi-tier problem.
